% Lecture slides for Guide 1, "Neural networks: a function with knobs".
% Written from content/01-neural-networks.mdx. Code, outputs, figures and notebook cell numbers
% are generated (tools/build_slide_assets.py, tools/export-figures.mjs); every other number
% must appear in the guide, which tools/check_slides.py checks. Build: node tools/build-slides.mjs
\providecommand\primeropts{}
\documentclass[aspectratio=169,11pt,t,\primeropts]{beamer}
\def\guidename{01-neural-networks}
\input{generated/01-neural-networks/meta.tex}
\input{generated/01-neural-networks/cells.tex}
\usetheme{primer}

\title{Neural networks}
\subtitle{A function with knobs}

\begin{document}

% ======================================================================== opening
\begin{frame}[plain,noframenumbering]
  \titlepage
\end{frame}

\begin{frame}{Where we are}
  \begin{center}\primerseriesmap\end{center}
  \vspace{2mm}
  \begin{itemize}
    \item Part 1 takes a language model apart, one piece at a time. This is the first guide: no
      earlier guide is assumed.
  \end{itemize}
\end{frame}

\begin{frame}{One recipe for the whole series}
  When you call an LLM (large language model) API, your text goes into a function and a reply
  comes out. That function's behavior is set by billions of numbers, and nobody wrote any of
  them by hand.
  \vspace{4mm}
  \begin{columns}[T]
    \column{0.31\textwidth}
    \begin{keyidea}[1 \textperiodcentered\ a function]
      with adjustable numbers
    \end{keyidea}
    \column{0.31\textwidth}
    \begin{keyidea}[2 \textperiodcentered\ a loss]
      one number that says how wrong the function currently is
    \end{keyidea}
    \column{0.31\textwidth}
    \begin{keyidea}[3 \textperiodcentered\ a procedure]
      that adjusts the first to shrink the second
    \end{keyidea}
  \end{columns}
  \vspace{4mm}
  Every later guide plugs different parts into this recipe. An LLM is trained with the same loop
  you'll write today.
\end{frame}

\begin{frame}{Learning objectives}
  By the end of this lecture you'll be able to:
  \begin{enumerate}
    \item \textbf{compute} a neuron's output on paper;
    \item \textbf{follow} an input through a two-layer network, shape by shape;
    \item \textbf{explain} how softmax turns scores into probabilities, and how cross-entropy
      turns them into one loss;
    \item \textbf{work out} a gradient with a calculator, by the nudge test and by
      backpropagation;
    \item \textbf{read} a short numpy training loop and \textbf{predict} how its error falls
      before you run it.
  \end{enumerate}
\end{frame}

% ======================================================================== section 1
\section{A network is a function with knobs}

\begin{frame}{Parameters and neurons}
  \begin{itemize}
    \item Each of those billions of numbers is a \term{parameter}: a knob whose setting changes
      what the function computes. We write $\theta$ for all of them together.
    \item A trained function of this kind is a model; in this guide the model is always a neural
      network. \textbf{Training} is the search for good settings.
    \item The smallest unit is a \term{neuron}:
      \begin{itemize}
        \item multiply each input by its own \term{weight},
        \item add up the results, then add one more number, the \term{bias},
        \item pass the sum through an \term{activation function}, a fixed rule applied to the
          single number.
      \end{itemize}
    \item Our activation throughout is \term{ReLU}, the rectified linear unit: $\max(0,\cdot)$.
      A positive number passes through, a negative one becomes zero
      (popularized for deep networks by Nair \& Hinton, 2010).
  \end{itemize}
\end{frame}

\begin{frame}{Worked example: one neuron}
  Input $x = [2, -1]$, weights $w = [1.5, 0.5]$, bias $b = 0.25$.
  \primerfigure[0.5\textheight]{1}
  \vspace{-2mm}
  \[
    \uncover<2->{\text{sum} = 2 \cdot 1.5 + (-1) \cdot 0.5 + 0.25}\uncover<3->{\, = 2.75}
    \uncover<4->{\qquad\longrightarrow\qquad \text{output} = \max(0,\, 2.75) = 2.75}
  \]
  \uncover<4->{\begin{center}\small The sum is positive, so ReLU passes it through unchanged.
    \quad{\color{primerMuted}(Figure 1 of the guide)}\end{center}}
\end{frame}

\begin{frame}{A second neuron, and the off switch}
  A second neuron reads the same input $x = [2, -1]$ with weights $[-1, 1]$ and bias $0$:
  \[
    \uncover<2->{(-1)\cdot 2 + 1\cdot(-1) + 0 = -2 - 1 = -3}
    \uncover<3->{\quad\longrightarrow\quad \max(0, -3) = 0}
  \]
  \uncover<3->{For this input the second neuron is \textbf{switched off}. The next section
  shows why that off switch is what makes stacking neurons worthwhile.}
  \vspace{4mm}
  \uncover<4->{%
  \begin{keyidea}[A name, not a mechanism]
    ``Neuron'' is borrowed from early mathematical models of brain cells. Nothing here depends
    on how brains work: a neuron in this series is a weighted sum followed by a simple fixed
    function, and nothing more.
  \end{keyidea}}
\end{frame}

\begin{frame}{Check your understanding}
  \checkq{Feed the input $x = [-1, 2]$ to both neurons (weights $[1.5, 0.5]$, bias $0.25$; and
    weights $[-1, 1]$, bias $0$). What does each output? Which numbers could training change?}
  {First: $(-1) \cdot 1.5 + 2 \cdot 0.5 + 0.25 = -0.25$, so ReLU outputs $0$. Second:
    $(-1)(-1) + 2 \cdot 1 + 0 = 3$, output $3$. Now the first neuron is off and the second on.
    Training changes the weights and biases; the input is data and ReLU is a fixed rule.}
\end{frame}

% ======================================================================== section 2
\section{Layers, and why the bend matters}

\begin{frame}{A layer is one matrix multiply}
  A \term{layer} is many neurons reading the same input. Give each neuron's weights their own
  \emph{column} of $W$ and line the biases up in a row $b$:
  \[ h = \sigma(xW + b) \]
  \begin{itemize}
    \item $x$: a row of $d_{\text{in}}$ numbers; $W$: $d_{\text{in}} \times d_{\text{out}}$;
      $b$ and $h$: rows of $d_{\text{out}}$ numbers, one per neuron.
    \item $\sigma$ is the activation, applied to each number separately.
    \item Inputs are rows and matrices multiply on the right, a convention the whole series keeps.
  \end{itemize}
  \pause
  Our two neurons are one layer:
  \[
    W = \begin{bmatrix} 1.5 & -1 \\ 0.5 & 1 \end{bmatrix},\quad b = [0.25, 0],\quad
    x = [2, -1] \;\Rightarrow\; h = [2.75, 0]
  \]

\end{frame}

\begin{frame}{Stacking layers: a network}
  Feed one layer's output into a second layer:
  \[
    h = \sigma\big(xW^{(1)} + b^{(1)}\big), \qquad \hat{y} = h W^{(2)} + b^{(2)}
  \]
  \begin{itemize}
    \item Superscripts number the layers; $\hat{y}$ is the network's output.
    \item The \textbf{output layer} is the same weighted sums with the activation left off, so
      its outputs can be any real numbers.
    \item A layer between input and output, like the first one here, is a \textbf{hidden layer}.
  \end{itemize}
\end{frame}

\begin{frame}{Why does the network need \texorpdfstring{$\sigma$}{the activation} at all?}
  Take $\sigma$ away and set the biases aside:
  \[
    \hat{y} = x W^{(1)} W^{(2)} \uncover<2->{= x\big(W^{(1)} W^{(2)}\big)}
  \]
  \uncover<3->{%
  Say the network reads 2 numbers, its first layer has 3 neurons and its second has 2:
  \[
    \underbrace{W^{(1)}}_{2 \times 3}\;\underbrace{W^{(2)}}_{3 \times 2}
    \;=\; \underbrace{\text{one matrix}}_{2 \times 2}
  \]}
  \uncover<4->{%
  \begin{itemize}
    \item Two layers without an activation compute exactly what one layer could.
    \item The same goes for a hundred of them.
    \item Putting the biases back only adds one combined bias to that single layer.
  \end{itemize}}
\end{frame}

\begin{frame}{ReLU breaks the collapse}
  \begin{columns}[c]
    \column{0.36\textwidth}
    \begin{center}
    \begin{tikzpicture}[scale=0.8]
      \draw[primerRuleStrong,->] (-2.2,0) -- (2.4,0) node[right,primerMuted,font=\scriptsize] {input};
      \draw[primerRuleStrong,->] (0,-0.4) -- (0,2.4) node[above,primerMuted,font=\scriptsize] {ReLU};
      \draw[primerAccent,line width=1.4pt] (-2,0) -- (0,0) -- (2.2,2.2);
      \fill[primerAccent] (0,0) circle (2.2pt);
      \node[font=\scriptsize,primerInk,anchor=north west] at (0.08,-0.05) {the bend};
    \end{tikzpicture}
    \end{center}
    \column{0.62\textwidth}
    \small
    \begin{itemize}
      \item A switched-off neuron contributes nothing, as if its column of $W$ were zeros: for
        any one input the layer acts like a single matrix plus a bias.
      \item<2-> \textbf{Which} matrix depends on the input:
        \begin{itemize}
          \item $x = [2, -1]$: sums $2.75$, $-3$; second off.
          \item $x = [-1, 2]$: sums $-0.25$, $3$; first off.
        \end{itemize}
      \item<3-> No fixed matrix behaves that way.
    \end{itemize}
  \end{columns}
  \uncover<4->{%
  \begin{keyidea}
    Each layer picks its matrix by which neurons are on: one straight-line piece per on/off
    pattern, stitched together. Far more than one matrix can represent.
  \end{keyidea}}
\end{frame}

\begin{frame}{How much can a network represent?}
  \begin{itemize}
    \item \textbf{Universal approximation theorem.} A network with a single hidden layer can
      approximate any continuous function on a closed, bounded region as closely as you like,
      given enough neurons (Cybenko, 1989; Hornik, Stinchcombe \& White, 1989).
    \item The original proofs covered S-shaped activations; Leshno et al.\ (1993) extended the
      result to any activation that isn't a polynomial, which includes ReLU.
  \end{itemize}
  \pause
  \begin{keyidea}[Read the theorem carefully]
    It only says that a good setting of the parameters \emph{exists}. It doesn't say how many
    neurons you need, and it doesn't say that training will find those parameters.
    ``Can represent'' is a much weaker promise than ``can learn.''
  \end{keyidea}
\end{frame}

\begin{frame}{Check your understanding}
  \checkq{A network reads 2 numbers, has a first layer of 3 neurons and a second of 2, and no
    activation anywhere. What single layer computes the same function?}
  {One layer with the $2 \times 2$ matrix $W^{(1)}W^{(2)}$ and one combined bias. Without the
    bend, depth adds nothing: the product of matrices is a matrix.}
\end{frame}

% ======================================================================== section 3
\section{The forward pass in shapes}

\begin{frame}{The task for the rest of the guide}
  Running an input through the layers to get an output is the \term{forward pass}. To follow
  one we need a problem:
  \begin{keyidea}[Task]
    Predict the next character of a text from the current character alone. The text uses 27
    characters: the space and the letters a to z.
  \end{keyidea}
  \vspace{2mm}
  Two declared simplifications, repaid later in the series:
  \begin{itemize}
    \item \textbf{One character of context.} An LLM takes into account all the earlier text it's
      given, up to a fixed limit (Guide 4).
    \item \textbf{One-hot input.} Each character becomes a row of 27 numbers, all zeros except a
      single 1 in its slot, a \term{one-hot vector}. LLMs feed in learned vectors instead
      (Guide 3).
  \end{itemize}
\end{frame}

\begin{frame}{Shapes, step by step}
  \vspace{-3mm}
  \primerfigure[0.69\textheight]{2}
  \vspace{-3mm}
  \begin{center}\footnotesize\color{primerMuted}
    16 hidden ReLU neurons: an arbitrary choice, plenty for this text \textperiodcentered\ Figure 2
    of the guide
  \end{center}
\end{frame}

\begin{frame}{Counting the parameters}
  Two weight matrices and two bias rows:
  \[
    \underbrace{27 \cdot 16}_{W^{(1)}} + \underbrace{16}_{b^{(1)}}
    + \underbrace{16 \cdot 27}_{W^{(2)}} + \underbrace{27}_{b^{(2)}}
    \uncover<2->{\;=\; 907}
  \]
  \uncover<3->{An LLM has billions as of October 2026; Guide 7 looks at what that scale costs.}
  \vspace{5mm}
  \uncover<4->{%
  \begin{keyidea}[Many inputs at once]
    Stack $n$ one-hot rows into an $n \times 27$ matrix, run the same two multiplies, and you get
    an $n \times 27$ matrix of scores, one row per input. Each row is computed independently,
    which is why the code later never loops over examples.
  \end{keyidea}}
\end{frame}

\begin{frame}{Check your understanding}
  \checkq{For a batch of $n$ characters, what are the shapes of the input, the hidden matrix
    and the scores? Which of them depend on $n$?}
  {Input $n \times 27$, hidden $n \times 16$, scores $n \times 27$: only the number of rows
    depends on $n$. The parameters ($27 \times 16$, $16$, $16 \times 27$, $27$) don't change,
    so the same 907 numbers serve every example.}
\end{frame}

% ======================================================================== section 4
\section{Measuring wrongness: softmax and cross-entropy}

\begin{frame}{From scores to probabilities: softmax}
  Scores that feed a softmax are \term{logits}, written $z$; they can be any real numbers. To
  read them as probabilities they must be positive and sum to 1.
  \[ p_i = \frac{e^{z_i}}{\sum_j e^{z_j}} \]
  \pause
  Worked example, logits $[2, 1, 0]$:
  \[ e^2 = 7.389,\ e^1 = 2.718,\ e^0 = 1 \ \Rightarrow\ \text{sum } 11.107
    \uncover<3->{\ \Rightarrow\ p = [0.665, 0.245, 0.090]} \]
  \vspace{-2mm}
  \uncover<4->{%
  \begin{keyidea}
    Logits one apart, each probability $e \approx 2.718$ times the next: softmax turns
    \emph{differences} between logits into \emph{ratios} between probabilities.
  \end{keyidea}}
\end{frame}

\begin{frame}{The network never outputs a character}
  \begin{itemize}
    \item It outputs a probability for every character: together, a \textbf{probability
      distribution} over the 27 options.
    \item Anything that looks like a single prediction is a \emph{choice} made afterwards from
      those 27 numbers.
  \end{itemize}
\end{frame}

\begin{frame}{One number for ``how wrong'': cross-entropy}
  \begin{columns}[T]
    \column{0.56\textwidth}
    A \term{loss} $\mathcal{L}$ is large when the network is badly wrong and small when it's
    nearly right: zero when the correct answer gets probability 1, growing without limit as that
    probability approaches zero.
    \[ \mathcal{L} = -\ln p_y \]
    \term{Cross-entropy}: the negative natural log of the probability given to the correct
    answer $y$.
    \begin{itemize}
      \item<2-> First option correct: $-\ln 0.665 = 0.408$.
      \item<3-> Probability 1: $-\ln 1 = 0$.
      \item<3-> Probability $0.01$: $4.6$.
    \end{itemize}
    \column{0.42\textwidth}
    \begin{center}
    \begin{tikzpicture}[x=4cm,y=0.62cm]
      \draw[primerRuleStrong,->] (0,0) -- (1.08,0) node[right,font=\scriptsize,primerMuted] {$p_y$};
      \draw[primerRuleStrong,->] (0,0) -- (0,5.3) node[above,font=\scriptsize,primerMuted] {$-\ln p_y$};
      \draw[primerAccent,line width=1.2pt,domain=0.0075:1,samples=120,smooth] plot (\x,{-ln(\x)});
      \node[font=\scriptsize,primerMuted,below] at (1,0) {1};
      \node[font=\scriptsize,primerMuted,below] at (0,0) {0};
      \only<2->{\fill[primerInk] (0.665,0.408) circle (1.6pt) node[above right,font=\scriptsize] {$0.408$};}
      \only<3->{\fill[primerInk] (0.01,4.6) circle (1.6pt) node[right,font=\scriptsize] {$4.6$};}
    \end{tikzpicture}
    \end{center}
  \end{columns}
\end{frame}

\begin{frame}{The baseline every run starts from}
  A network with no information spreads its probability evenly: $1/27$ each, so
  \[ \mathcal{L} = -\ln(1/27) = \ln 27 \approx 3.30 \]
  \begin{itemize}
    \item Our toy network starts almost exactly there: small random weights give logits close
      to zero, and equal logits mean equal probabilities.
    \item Any loss below $3.30$ means the network has learned something about the text.
    \item An \textbf{example} is one current character paired with the character that follows
      it. The loss we train on is the average over all examples.
    \item An LLM's last step is this same softmax, run over its whole \textbf{vocabulary}, the set
      of tokens it can output. For now think of a token as roughly a word: Guide 2 refines
      that, and Guide 6 builds on the softmax.
  \end{itemize}
\end{frame}

\begin{frame}{Check your understanding}
  \checkq{The network gives probabilities $[0.665, 0.245, 0.090]$ and the first option is
    correct. What is the loss? And why is $\ln 27$ the natural reference value for our task?}
  {$-\ln 0.665 = 0.408$. A network that knows nothing gives each of the 27 characters
    probability $1/27$, whose loss is $\ln 27 \approx 3.30$; anything lower is evidence of
    learning.}
\end{frame}

% ======================================================================== section 5
\section{Learning is walking downhill, precisely}

\begin{frame}{One knob, by hand: the nudge test}
  Which way should each of the 907 knobs move? Start with one knob and
  $\mathcal{L}(\theta) = (\theta - 3)^2$.
  \begin{align*}
    \theta = 0 &: \quad \mathcal{L} = 9\\
    \uncover<2->{\theta = 0.01 &: \quad \mathcal{L} = 2.99^2 = 8.9401}
  \end{align*}
  \uncover<3->{A drop of about $0.06$ for a nudge of $0.01$: the loss changes about $-6$ times as
  fast as the knob.}
  \vspace{3mm}
  \uncover<4->{%
  \begin{itemize}
    \item That ratio is the slope. The \term{gradient} is that slope, measured with a nudge small
      enough that the ratio stops depending on its size.
    \item For this loss the slope at any $\theta$ is $2(\theta - 3)$.
    \item Many knobs: one slope per knob, each measured with the others fixed, written
      $\nabla_\theta \mathcal{L}$.
  \end{itemize}}
\end{frame}

\begin{frame}{Gradient descent}
  Move each knob against the sign of its slope, further when the slope is steeper:
  \[ \theta \leftarrow \theta - \eta \, \nabla_\theta \mathcal{L} \]
  \term{Gradient descent}; the multiplier $\eta$ is the \term{learning rate}.
  \pause
  \vspace{1mm}
  With $\eta = 0.1$, starting at $\theta = 0$:
  \begin{align*}
    \text{step 1:}\quad & 0 - 0.1 \cdot (-6) = 0.6\\
    \uncover<3->{\text{step 2:}\quad & \text{slope } 2(0.6 - 3) = -4.8 \;\Rightarrow\; \theta = 1.08}
  \end{align*}
  \uncover<4->{Every step ends closer to $3$, and the steps shrink as the slope flattens.}
  \uncover<5->{%
  \begin{keyidea}[No guarantee in general]
    Gradient descent reaches the best setting here only because this loss is bowl-shaped. For a
    network, a small enough step lowers the loss a little, but the walk can end anywhere the
    slope has gone flat.
  \end{keyidea}}
\end{frame}

\begin{frame}{Two knobs: walking downhill, literally}
  \begin{columns}[T]
    \column{0.29\textwidth}
    \footnotesize
    Loss $\theta_1^2 + 4\theta_2^2$, $\eta = 0.1$; slopes $2\theta_1$ and $8\theta_2$.
    \begin{itemize}
      \item Start $(4, 1.5)$: slopes $8$ and $12$.
      \item First step: $(4 - 0.8,\; 1.5 - 1.2)$ $= (3.2, 0.3)$.
      \item Loss $25 \to 10.6$.
    \end{itemize}
    Dashed: the gradient, pointing uphill, crossing contours at a right angle. Solid: the
    update, the opposite way.
    \vspace{2mm}

    {\color{primerMuted}Figure 3 of the guide}
    \column{0.69\textwidth}
    \vspace{-5mm}
    \primerfigure[0.76\textheight]{3}
  \end{columns}
\end{frame}

\begin{frame}{The length of the stride}
  Read as a direction, the list of slopes points where the loss rises fastest (used without
  proof). The update steps the opposite way, and $\eta$ sets the length of the stride.
  \vspace{2mm}
  \begin{columns}[T]
    \column{0.31\textwidth}
    \begin{keyidea}[Too long]
      $\eta = 1.1$: the first step from $0$ lands on $6.6$, further from $3$ than where it
      started.
    \end{keyidea}
    \column{0.31\textwidth}
    \begin{keyidea}[Too short]
      $\eta = 0.001$: the first step only reaches $0.006$. It crawls.
    \end{keyidea}
    \column{0.31\textwidth}
    \begin{keyidea}[Bounces forever]
      $\eta = 1$ on this loss bounces between $0$ and $6$ forever.
    \end{keyidea}
  \end{columns}
  \vspace{3mm}

  {\small\color{primerMuted}Declared simplification: every step here measures the loss on the
  whole dataset. Real training uses a small random sample per step and more elaborate update
  rules (Guide 7).}
\end{frame}

\begin{frame}{Check your understanding}
  \checkq{On $\mathcal{L}(\theta) = (\theta - 3)^2$: what is the gradient at $\theta = 0.6$?
    With $\eta = 0.1$, where does the next step land, and why are the steps getting shorter?}
  {$2(0.6 - 3) = -4.8$, so $\theta \leftarrow 0.6 - 0.1 \cdot (-4.8) = 1.08$. Each step is
    $\eta$ times the slope, and the slope shrinks as $\theta$ approaches $3$: $-6$, then
    $-4.8$.}
\end{frame}

% ======================================================================== section 6
\section{Backpropagation: the chain rule, run backward}

\begin{frame}{Why we need something better than nudging}
  \begin{itemize}
    \item The nudge test gives one slope per extra forward pass: 907 extra passes for every step.
    \item \term{Backpropagation} gets all of them for about the cost of one more pass.
    \item It rests on the \term{chain rule}: rates along a chain multiply.
  \end{itemize}
  \vspace{3mm}
  \begin{center}
  \begin{tikzpicture}[node distance=2.6cm,
      q/.style={draw=primerRuleStrong,rounded corners=2pt,minimum width=1.6cm,minimum height=0.8cm,font=\small}]
    \node[q] (a) {first};
    \node[q,right=of a] (b) {second};
    \node[q,right=of b] (c) {third};
    \draw[-{Stealth},primerInk] (a) -- node[above,font=\small] {rate 2} (b);
    \draw[-{Stealth},primerInk] (b) -- node[above,font=\small] {rate 3} (c);
    \draw[-{Stealth},primerAccent,line width=1pt] (a.south) to[bend right=18] node[below,font=\small] {rate $2 \times 3 = 6$} (c.south);
  \end{tikzpicture}
  \end{center}
\end{frame}

\begin{frame}{A network small enough for paper: forward}
  Input $x = 1$; weights $w_1 = 2$, $w_2 = 0.5$; correct answer $y = 2$.
  \vspace{-2mm}
  \[ h = \text{ReLU}(w_1 x), \qquad \hat{y} = w_2 h, \qquad
    \mathcal{L} = \tfrac{1}{2}(\hat{y} - y)^2 \]
  \vspace{-4mm}
  \primerfigonly[0.56\textheight]{4}{1}
  \vspace{-3mm}
  \[
    \uncover<2->{h = \text{ReLU}(2 \cdot 1) = 2}\qquad
    \uncover<3->{\hat{y} = 0.5 \cdot 2 = 1}\qquad
    \uncover<4->{\mathcal{L} = \tfrac{1}{2}(1 - 2)^2 = 0.5}
  \]
  {\footnotesize\color{primerMuted}Figure 4 of the guide, forward view.\par}
\end{frame}

\begin{frame}{The backward pass, one rate at a time}
  Start at the loss and work toward the input. $\partial\mathcal{L}/\partial w$ is the slope of
  the loss with respect to $w$, everything else fixed. (Squared error only keeps the arithmetic
  small; the steps are the same for cross-entropy.)
  \vspace{2mm}
  \begin{enumerate}
    \item<1-> \textbf{Output.} The slope of $\tfrac{1}{2}(\hat{y} - y)^2$ is $\hat{y} - y$ (the
      $\tfrac{1}{2}$ cancels the 2): $\partial\mathcal{L}/\partial\hat{y} = 1 - 2 = -1$.
    \item<2-> \textbf{Second weight.} Since $\hat{y} = w_2 h$, nudging $w_2$ moves $\hat{y}$ at rate $h = 2$:
      $\partial\mathcal{L}/\partial w_2 = -1 \cdot 2 = -2$.
    \item<3-> \textbf{Hidden value.} $\hat{y}$ moves with $h$ at rate $w_2 = 0.5$:
      $\partial\mathcal{L}/\partial h = -1 \cdot 0.5 = -0.5$.
    \item<4-> \textbf{ReLU.} Its input, $w_1 x = 2$, is positive: rate 1, so the slope is still
      $-0.5$.
    \item<5-> \textbf{First weight.} Nudging $w_1$ moves ReLU's input, $w_1 x$, at rate $x = 1$:
      $\partial\mathcal{L}/\partial w_1 = -0.5 \cdot 1 = -0.5$.
  \end{enumerate}
\end{frame}

\begin{frame}{The backward pass, drawn}
  \vspace{-2mm}
  \primerfigonly[0.68\textheight]{4}{2}
  \vspace{-3mm}
  \begin{center}\small
    Gradients flow right to left: $-1$ at the output, $-2$ for $w_2$, $-0.5$ at $h$, $-0.5$ for
    $w_1$. {\color{primerMuted}(Figure 4 of the guide, backward view)}
  \end{center}
\end{frame}

\begin{frame}{One step of gradient descent, and the payoff}
  With $\eta = 0.1$:
  \begin{align*}
    w_2 &= 0.5 - 0.1 \cdot (-2) = 0.7 & w_1 &= 2 - 0.1 \cdot (-0.5) = 2.05
  \end{align*}
  \pause
  Forward again:
  \[
    h = \text{ReLU}(2.05 \cdot 1) = 2.05, \quad \hat{y} = 0.7 \cdot 2.05 = 1.435, \quad
    \mathcal{L} = \tfrac{1}{2}(1.435 - 2)^2 \approx 0.160
  \]
  \pause
  \begin{keyidea}
    The loss fell from $0.5$ to $0.160$ in one step, using the two gradients backpropagation
    gave us.
  \end{keyidea}
\end{frame}

\begin{frame}{Why one backward pass costs about one forward pass}
  The backward pass reused two things:
  \begin{itemize}
    \item \textbf{Forward values.} Step 2 needed $h$, which the forward pass had already
      computed, so a real implementation keeps the forward values around.
    \item \textbf{Shared slopes.} Step 4 started from $\partial\mathcal{L}/\partial h$, computed
      once in step 3, and step 5 built on step 4. In a deeper network every weight earlier than
      $h$ would start from that same number.
  \end{itemize}
  \vspace{2mm}
  \begin{keyidea}
    Each intermediate slope is computed once and shared by everything upstream of it. That
    sharing is why one backward pass costs about as much as one forward pass, however many
    parameters there are.
  \end{keyidea}
\end{frame}

\begin{frame}{Two results that carry over to the 907-parameter network}
  \begin{enumerate}
    \item \textbf{ReLU gates the gradient.} Its rate is 1 where its input was positive and 0
      where it wasn't: a switched-off neuron sends no gradient back, and its incoming weights
      don't move for that example.
    \item<2-> \textbf{Softmax + cross-entropy: predicted minus correct.} The slope of the loss
      with respect to each logit is the predicted probabilities minus the correct answer's
      one-hot vector.
  \end{enumerate}
  \uncover<3->{%
  \[
    p = [0.665,\; 0.245,\; 0.090],\ \text{first correct} \;\Rightarrow\;
    \frac{\partial\mathcal{L}}{\partial z} = [-0.335,\; 0.245,\; 0.090]
  \]
  The updates move to raise the first logit and lower the other two.}
\end{frame}

\begin{frame}{Going deeper (optional): backpropagation for the matrix network}
  For $n$ inputs as rows: $X$ ($n \times 27$), hidden $H$ ($n \times 16$), probabilities $P$
  ($n \times 27$), one-hot answers $Y$ ($n \times 27$). Every gradient has the shape of the thing
  it is a gradient for.
  \vspace{2mm}
  \begin{center}\small
  \begin{tabular}{@{}lL{0.62\textwidth}l@{}}
    \toprule
    \textbf{Gradient on} & \textbf{Computed as} & \textbf{Shape}\\
    \midrule
    logits & $(P - Y)/n$ (division: the loss is a mean) & $n \times 27$\\
    $W^{(2)}$ & the transpose of $H$ times the logit gradient & $16 \times 27$\\
    $b^{(2)}$ & column sums of the logit gradient & $1 \times 27$\\
    $H$ & the logit gradient times the transpose of $W^{(2)}$, zeroed where ReLU's output was zero & $n \times 16$\\
    $W^{(1)}$ & the transpose of $X$ times that & $27 \times 16$\\
    $b^{(1)}$ & its column sums & $1 \times 16$\\
    \bottomrule
  \end{tabular}
  \end{center}
  \vspace{2mm}
  {\small These are the lines of \texttt{backward}, cell In [\primercell{235}].}
\end{frame}

\begin{frame}{Going deeper (optional): where ``predicted minus correct'' comes from}
  Two facts from calculus: the slope of $\ln u$ is $1/u$; the slope of $e^u$ is $e^u$.
  \pause
  \[
    \mathcal{L} = -\ln p_y = -z_y + \ln \sum_j e^{z_j}
  \]
  \pause
  Nudge $z_i$:
  \begin{itemize}
    \item the second term changes at rate $\dfrac{e^{z_i}}{\sum_j e^{z_j}} = p_i$;
    \item the first term contributes $-1$ when $i$ is the correct answer, nothing otherwise.
  \end{itemize}
  \pause
  \[
    \frac{\partial\mathcal{L}}{\partial z_i} =
    \begin{cases} p_i - 1 & \text{for the correct answer, } i = y\\ p_i & \text{for every other option} \end{cases}
  \]
  \vspace{-1mm}
  \begin{center}which is $p$ minus the one-hot vector.\end{center}
  {\small\color{primerMuted}Backpropagation became standard after Rumelhart, Hinton \& Williams
  (1986) showed it learning useful hidden layers; the method is older (e.g.\ Linnainmaa, 1976).}
\end{frame}


\begin{frame}{Check your understanding}
  \checkq{In the paper network ($x = 1$, $w_1 = 2$, $w_2 = 0.5$, $y = 2$), which rates multiply
    to give $\partial\mathcal{L}/\partial w_2$, and what is it? And if ReLU's input had been
    negative, what would $\partial\mathcal{L}/\partial w_1$ have been?}
  {$\partial\mathcal{L}/\partial\hat{y} = -1$ times the rate $h = 2$: $-2$. With a negative
    input ReLU's rate is 0, so no gradient reaches $w_1$: it would be $0$, and $w_1$ wouldn't
    move for that example.}
\end{frame}

% ======================================================================== section 7
\section{The whole loop in numpy}

\begin{frame}{The data: characters to training examples}
  Each row of \texttt{X} is a one-hot current character; the matching entry of \texttt{y} is
  the index of the character that follows it.
  \vspace{1mm}
  \primercode{179}
  \primeroutput{179}
  \cellnote{179}
\end{frame}

\begin{frame}{The parameters: small random numbers, fixed seed}
  \begin{columns}[T]
    \column{0.55\textwidth}
    \primercode{199}
    \primeroutput{199}
    \cellnote{199}
    \column{0.42\textwidth}
    \small
    \begin{itemize}
      \item They can't all start at zero: every hidden neuron would compute the same thing and
        receive the same gradient, and could never become different.
      \item The fixed seed makes your run match the guide's (NumPy 2.2.6, pinned as of October
        2026).
      \item 907, as counted by hand.
    \end{itemize}
  \end{columns}
\end{frame}

\begin{frame}{Forward pass and loss, for all $n$ rows at once}
  \primercode{214}
  \primeroutput{214}
  \cellnote{214}
\end{frame}

\begin{frame}{Reading the forward pass}
  \begin{itemize}
    \item The loss prints $3.296$: the $\ln 27$ of a uniform guess, as predicted.
    \item \texttt{h} and \texttt{z} are the two layer equations, for every row at once.
    \item One line subtracts each row's largest logit from the whole row. Softmax depends only
      on the differences between logits, so the probabilities don't change, and \texttt{exp} is
      never handed a large positive number.
    \item \texttt{loss\_of} picks each row's probability for its correct next character and
      averages $-\ln$ of them: mean cross-entropy.
  \end{itemize}
\end{frame}

\begin{frame}{The backward pass, one line per gradient}
  \primercode{235}
  \cellnote{235}
\end{frame}

\begin{frame}{From the paper steps to \texttt{backward}}
  Each line is one of the paper steps, done for every weight and every example at once.
  \vspace{2mm}
  \begin{center}\small
  \begin{tabular}{@{}L{0.27\textwidth}L{0.62\textwidth}@{}}
    \toprule
    \textbf{Line} & \textbf{Paper step and reason}\\
    \midrule
    \texttt{dz} & predicted minus correct on the logits, divided by $n$: the loss is a mean\\
    \texttt{dW2 = h.T @ dz} & the second-weight step (hidden value times logit slope), added up
      over examples: slopes add because the loss averages per-example losses\\
    \texttt{db2} & a bias moves its logit one-for-one: the column sum\\
    \texttt{dh = dz @ W2.T} & the hidden-value step\\
    \texttt{dh[h <= 0] = 0} & ReLU: no gradient where the neuron was off\\
    \texttt{dW1}, \texttt{db1} & the same pattern, one layer down\\
    \bottomrule
  \end{tabular}
  \end{center}
\end{frame}

\begin{frame}{Training: predict before you run}
  \begin{columns}[T]
    \column{0.52\textwidth}
    \renewcommand\codesize{\scriptsize}
    \primercode{253}
    \cellnote{253}
    \vspace{1mm}
    {\small $\eta = 1$ would bounce forever on the one-knob loss. How long is too long depends on
    how sharply the particular loss curves, so a learning rate is found by trying.\par}
    \column{0.45\textwidth}
    \begin{keyidea}[Predict first]
      Start $3.296$, $\eta = 1$, 2,000 steps. How fast does the loss fall, and does it reach
      zero?
    \end{keyidea}
    \only<2-| handout:1>{%
      \primeroutput{253}
      \vspace{1mm}
      {\small Fast at first: $1.176$ after 100 steps. Then it levels off just below $0.89$;
      another thousand steps buy only $0.003$.\par}}
  \end{columns}
\end{frame}

\begin{frame}{Why it stops: the text sets a floor}
  In the text, ``t'' is followed by another character 19 times: 11 times by ``h'', 7 times by a
  space and once by ``o'': 58\%, 37\% and 5\%.
  \vspace{1mm}
  \begin{itemize}
    \item A function of the current character alone can't tell those cases apart, so it gives
      one set of probabilities for all of them.
    \item<2-> Average cross-entropy is lowest when the probabilities equal the true frequencies.
  \end{itemize}
  \uncover<3->{%
  \begin{align*}
    \text{put } 0.90, 0.05, 0.05:&\quad
      \tfrac{11}{19}(-\ln 0.90) + \tfrac{7}{19}(-\ln 0.05) + \tfrac{1}{19}(-\ln 0.05) \approx 1.32\\
    \text{use } \tfrac{11}{19}, \tfrac{7}{19}, \tfrac{1}{19}:&\quad \approx 0.84
      \quad\text{(the floor for the 19 examples after ``t'')}
  \end{align*}}
\end{frame}

\begin{frame}{The best possible table, measured}
  \primercode{279}
  \primeroutput{279}
  \cellnote{279}
  \vspace{1mm}
  \begin{itemize}
    \item No function of the current character can get below $0.884$ on this text; the network
      reached $0.887$.
    \item<2-> A claim about the 127 training examples only; real training is judged on unseen
      text (Guide 7). Doing better needs more context: attention, in Guide 4.
  \end{itemize}
\end{frame}

\begin{frame}{Inference: probabilities close to the text's frequencies}
  \primercode{295}
  \primeroutput{295}
  \cellnote{295}
\end{frame}

\begin{frame}{Reading the inference output}
  \begin{itemize}
    \item After ``t'' the most likely next character is ``h''; after ``e'', a space. Both sit close
      to the text's frequencies.
    \item The network still puts $0.32$ on a space after ``t'': naming ``h'' as its prediction is a
      choice made from those probabilities.
    \item The loop was \term{training}: every pass changed the parameters. This block was
      \term{inference}: forward pass only, parameters exactly as training left them.
  \end{itemize}
\end{frame}

\begin{frame}{Check your understanding}
  \checkq{Someone suggests training for many more steps to push the loss on this text to $0.5$.
    What do you tell them?}
  {It can't happen with this input. The best any function of the current character can do on
    this text is $0.884$, and the network is already at $0.887$. A lower loss needs more
    context, not more steps.}
\end{frame}

% ======================================================================== section 8
\section{What carries forward to LLMs}

\begin{frame}{From the toy to an LLM}
  Each piece has a counterpart you'll meet later (sizes as of October 2026):
  \vspace{2mm}
  \begin{center}\small
  \begin{tabular}{@{}L{0.34\textwidth}L{0.44\textwidth}l@{}}
    \toprule
    \textbf{Toy network} & \textbf{LLM} & \textbf{Guide}\\
    \midrule
    one-hot input & a learned vector for each token & 2, 3\\
    the single current character & all the earlier text, through attention & 4\\
    two layers & dozens of transformer blocks & 5\\
    softmax over 27 characters & softmax over tens or hundreds of thousands of tokens & 6\\
    907 parameters, one sentence & billions of parameters, trillions of tokens & 7\\
    \bottomrule
  \end{tabular}
  \end{center}
\end{frame}

\begin{frame}{What stays the same: the recipe}
  \begin{columns}[T]
    \column{0.23\textwidth}
    \begin{keyidea}[function]
      with adjustable numbers
    \end{keyidea}
    \column{0.23\textwidth}
    \begin{keyidea}[loss]
      says how wrong it is
    \end{keyidea}
    \column{0.23\textwidth}
    \begin{keyidea}[backpropagation]
      gives the gradient of that loss
    \end{keyidea}
    \column{0.23\textwidth}
    \begin{keyidea}[descent step]
      moves every parameter a little way against it
    \end{keyidea}
  \end{columns}
  \vspace{5mm}
  \begin{center}\large An LLM is trained with that same loop.\end{center}
\end{frame}

% ======================================================================== closing
\section*{Wrap-up}

\begin{frame}{What we simplified, and where it's repaid}
  \begin{center}\small
  \begin{tabular}{@{}L{0.36\textwidth}L{0.42\textwidth}l@{}}
    \toprule
    \textbf{In this guide} & \textbf{In an LLM} & \textbf{Repaid}\\
    \midrule
    \simplification{one character of context}{all earlier text, up to a fixed limit}{Guide 4}
    \simplification{one-hot inputs}{learned vectors}{Guide 3}
    \simplification{each step on the whole dataset, plain gradient descent}{small random samples, more elaborate update rules}{Guide 7}
    \simplification{a token is ``roughly a word''}{refined in Guide 2}{Guide 2}
    \simplification{loss on the training examples}{loss on text the model never saw}{Guide 7}
    \bottomrule
  \end{tabular}
  \end{center}
  \vspace{2mm}
  {\small\color{primerMuted}The paper backpropagation example used squared error only for small
  arithmetic; language models train on cross-entropy.}
\end{frame}

\begin{frame}{Summary}
  \begin{enumerate}\small
    \item A neuron is a weighted sum plus a bias through a fixed activation; its weights and bias
      are parameters.
    \item A layer is one matrix multiply. Without the bend, stacked layers collapse into one;
      ReLU makes the matrix depend on the input.
    \item The forward pass maps a $1 \times 27$ one-hot row to 27 logits through 907 parameters,
      and handles $n$ rows at once.
    \item Softmax turns logit differences into probability ratios; cross-entropy $-\ln p_y$
      scores them, starting near $\ln 27 \approx 3.30$.
    \item Gradient descent moves every knob against its slope, scaled by $\eta$: too long a stride
      overshoots, too short crawls.
    \item Backpropagation multiplies rates backward, sharing each slope, so all gradients cost
      about one extra pass. On the logits it gives predicted minus correct.
    \item The numpy loop reaches $0.887$ against a floor of $0.884$: the limit is the
      one-character context, not the training.
  \end{enumerate}
\end{frame}

\begin{frame}{Exercises: try it in the notebook}
  \small
  \begin{enumerate}
    \item \textbf{On paper.} \free{For $x = [1, 3]$, $w = [0.5, -1]$, $b = 2$,} compute the
      neuron's sum and its ReLU output. Then find an input that switches it on.
    \item \textbf{Gradient check} (cell In [\primercell{235}]). Pick one entry of \texttt{W2},
      nudge it by \free{$10^{-5}$}, and compare the change in loss per unit nudge with the
      matching entry of \texttt{dW2}. Explain any small difference.
    \item \textbf{Learning rate} (cell In [\primercell{253}]). Rerun training with
      \free{$\eta = 0.1$ and $\eta = 3$}. Before running, predict which one crawls and which
      one might overshoot.
    \item \textbf{Hidden size} (cells In [\primercell{199}]--[\primercell{253}]). Change 16 hidden
      neurons to \free{4, then 64}. Predict the parameter count from the formula, then check
      whether the loss still approaches the floor of $0.884$.
    \item \textbf{Symmetry} (cell In [\primercell{199}]). Fill \texttt{W1} and \texttt{W2} with
      the same small positive number, \free{say 0.01}, so every hidden neuron is on and identical.
      Train, then compare the columns of \texttt{W1}. (Why not all zeros? Try it, and look at
      \texttt{h}.)
  \end{enumerate}
\end{frame}

\begin{frame}{Further reading}
  \small
  \begin{itemize}
    \item V.\ Nair \& G.\ Hinton (2010). Rectified linear units improve restricted Boltzmann
      machines. \emph{ICML}. \ (ReLU)
    \item G.\ Cybenko (1989). Approximation by superpositions of a sigmoidal function.
      \emph{Mathematics of Control, Signals and Systems}.
    \item K.\ Hornik, M.\ Stinchcombe \& H.\ White (1989). Multilayer feedforward networks are
      universal approximators. \emph{Neural Networks}.
    \item M.\ Leshno, V.\ Lin, A.\ Pinkus \& S.\ Schocken (1993). Multilayer feedforward networks
      with a nonpolynomial activation function can approximate any function. \emph{Neural
      Networks}.
    \item D.\ Rumelhart, G.\ Hinton \& R.\ Williams (1986). Learning representations by
      back-propagating errors. \emph{Nature}.
    \item S.\ Linnainmaa (1976). Taylor expansion of the accumulated rounding error. \emph{BIT}.
  \end{itemize}
  {\small\color{primerMuted}Links for each are in the \href{\guideurl}{guide}.}
\end{frame}

\begin{frame}{Next time: Guide 2, tokenization}
  \begin{itemize}
    \item Today's network read characters, one-hot. An LLM's softmax runs over its vocabulary,
      the set of tokens it can output, and we called a token ``roughly a word''.
    \item Guide 2 shows what a token really is.
  \end{itemize}
  \vspace{4mm}
  \begin{keyidea}[Before next time]
    Run the notebook from top to bottom and check that every printed number matches the slides.
  \end{keyidea}
\end{frame}


\end{document}
